\documentclass[aspectratio=169,11pt]{beamer}

\usepackage{amsmath,amssymb}
\usetheme{Madrid}

\title{First-Order Logic}
\author{}
\date{}

\begin{document}

%------------------------------------------------
%------------------------------------------------
%------------------------------------------------
%------------------------------------------------
\section{First Order Logic}
\begin{frame}{Question 1}
\textbf{Domain:} Integers.

\textbf{Predicate:} $L(x,y)$: ``$x<y$.''

\begin{enumerate}
    \item Translate:
    ``For every integer $y$, there exists an integer $x$ such that $x<y$.''

    \item Translate:
    ``There exists an integer $x$ such that for every integer $y$,
    $x<y$.''

    \item Are these two statements equivalent? Explain using
    standard arithmetic.
\end{enumerate}

\begin{block}{Answer}
\begin{enumerate}
    \item $\forall y\,\exists x\,L(x,y)$
    \item $\exists x\,\forall y\,L(x,y)$
    \item Not equivalent.
\end{enumerate}

$\forall y\,\exists x\,L(x,y)$ is \textbf{True}: for any $y$, choose $x=y-1$.

$\exists x\,\forall y\,L(x,y)$ is \textbf{False}: no single integer is smaller than every integer.

\textbf{Conclusion:} Different quantifier orders can have different truth values.
\end{block}

\end{frame}
%------------------------------------------------
\begin{frame}{Question 2}

\textbf{Domain:} Natural numbers.

\textbf{Function:} $s(x)$: successor of $x$ (i.e., $x+1$).

\textbf{Predicate:} $E(x)$: ``$x$ is even.''


\textbf{Express the following in first-order logic:}

\begin{enumerate}
    \item Zero is even. (Assume constant $0$.)
    
    \item If $x$ is even, then $s(s(x))$ is even.
    
    \item There exists an even number.
\end{enumerate}

\begin{block}{Answer}
\begin{enumerate}
    \item $ E(0)$
    \item 
    $\forall x\,(E(x)\to E(s(s(x))))$
    \item 
    $\exists x\,E(x)$
    
\end{enumerate}
\end{block}

\end{frame}
%------------------------------------------------
\begin{frame}{Question 3}

\textbf{Negate and simplify }

\begin{enumerate}
    \item
    $\forall x\,(P(x)\to Q(x))$
    \item
    $\exists x\,(P(x)\land\forall y\,R(x,y))$
\end{enumerate}
\begin{block}{Answer}

\textbf{1.}

\[
\begin{aligned}
\neg\forall x\,(P(x)\to Q(x))
&\equiv \exists x\,\neg(P(x)\to Q(x))\\
&\equiv \exists x\,(P(x)\land\neg Q(x)).
\end{aligned}
\]


\textbf{2.}

\[
\begin{aligned}
\neg\exists x\,(P(x)\land\forall y\,R(x,y))
&\equiv \forall x\,\neg(P(x)\land\forall y\,R(x,y))\\
&\equiv \forall x\,(\neg P(x)\lor\neg\forall y\,R(x,y))\\
&\equiv \forall x\,(\neg P(x)\lor\exists y\,\neg R(x,y)).
\end{aligned}
\]

\end{block}
\end{frame}

\begin{frame}{Question 4}
Given
$$
\varphi=\forall x\,\exists y\,P(x,y),
$$
which of the following is equivalent to $\neg\varphi$?

\begin{enumerate}
    \item[(A)] $\exists x\,\forall y\,\neg P(x,y)$
    \item[(B)] $\forall x\,\exists y\,\neg P(x,y)$
    \item[(C)] $\exists x\,\neg(\exists y\,P(x,y))$
    \item[(D)] $\neg\forall x\,\exists y\,P(x,y)$
\end{enumerate}
\begin{block}{Answer: (A), (C), and (D)}
$$
\neg\forall x\,\exists y\,P(x,y)
\equiv
\exists x\,\neg\exists y\,P(x,y)
\equiv
\exists x\,\forall y\,\neg P(x,y)
$$

Hence, the correct answers are
$$
\text{(A), (C), and (D)}.
$$
\end{block}
\end{frame}

%------------------------------------------------
\begin{frame}{Question 5}
Which one of the following first-order logic formulas is \textbf{valid}
(true in every interpretation with a non-empty domain)?

\begin{enumerate}
    \item[(A)]
    $$\forall x\,(P(x)\lor Q(x))
    \to
    (\forall x\,P(x)\lor\forall x\,Q(x))$$

    \item[(B)]
    $$\exists x\,(P(x)\land Q(x))
    \leftrightarrow
    (\exists x\,P(x)\land\exists x\,Q(x))$$

    \item[(C)]
    $$\forall x\,(P(x)\to Q(x))
    \to
    (\forall x\,P(x)\to\forall x\,Q(x))$$

    \item[(D)]
    $$\exists x\,(P(x)\to Q(x))
    \to
    (\exists x\,P(x)\to\exists x\,Q(x))$$
\end{enumerate}
\end{frame}

\begin{frame}
\begin{block}{Answer: (C)}
\textbf{(A):} Not valid. $P$ and $Q$ can hold for different elements.

\textbf{(B):} Not valid. $P$ and $Q$ may each hold for some element,
but not for the same element.

\textbf{(C):} Valid. If
$$
\forall x(P(x)\to Q(x))
\quad\text{and}\quad
\forall xP(x),
$$
then for every $x$, $P(x)$ is true. Hence $Q(x)$ is also true:
$$
\forall xQ(x).
$$

\textbf{(D):} Not valid. The witness for
$$
\exists x(P(x)\to Q(x))
$$
may be an element where $P(x)$ is false, while $\exists xP(x)$
is true and $\exists xQ(x)$ is false.

Therefore,
$$
\boxed{\text{(C)}}
$$
is valid.
\end{block}
\end{frame}


%------------------------------------------------
\section{Modus Ponens and Modus Tollens}
\begin{frame}{Question 1}
Given:

\begin{itemize}
    \item Premise 1: If a number is divisible by 4, then it is even.
    \item Premise 2: The number is even.
    \item Conclusion: Therefore, the number is divisible by 4.
\end{itemize}

Which statement is correct?

\begin{enumerate}
    \item[(A)] The argument is valid by Modus Ponens.
    \item[(B)] The argument is valid by Modus Tollens.
    \item[(C)] The argument is invalid (fallacy of affirming the consequent).
    \item[(D)] The argument is invalid (fallacy of denying the antecedent).
\end{enumerate}
\end{frame}

\begin{frame}
\begin{block}{Answer: (C)}
Let
$$
p:\text{number is divisible by 4},\qquad
q:\text{number is even}.
$$

The argument has the form
$$
p\to q,\qquad q,\qquad\therefore p.
$$

This is \textbf{affirming the consequent}, which is invalid.

For example,
$$
6\text{ is even, but }6\text{ is not divisible by }4.
$$

Hence,
$$
\boxed{\text{(C)}}
$$
\end{block}
\end{frame}
%------------------------------------------------
\begin{frame}{Question 2}
Given the premises:

\begin{enumerate}
    \item If the server is down, then users cannot log in.
    \item Users can log in.
\end{enumerate}
$$
p:\text{server is down},\qquad
q:\text{users cannot log in}.
$$
$$
p\to q,\qquad \neg q.
$$

Which conclusion is logically valid?

\begin{enumerate}
    \item[(A)] $p$
    \item[(B)] $\neg p$
    \item[(C)] $q\to p$
    \item[(D)] Nothing definite can be concluded.
\end{enumerate}
\end{frame}

\begin{frame}
\begin{block}{Answer: (B) only}
The premises are
$$
p\to q,\qquad \neg q.
$$

By \textbf{Modus Tollens},
$$
p\to q,\quad \neg q
\;\Rightarrow\;
\neg p.
$$

Therefore, the server is not down.

Option (C) is only the converse of the given implication, and
option (D) is false because $\neg p$ can be concluded.

Hence,
$$
\boxed{\text{(B)}}
$$
\end{block}
\end{frame}
\begin{frame}{Question 3}
\textbf{Argument:}

\begin{itemize}
    \item Premise 1: \hspace{1cm}
    \item Premise 2: The alarm is not ringing.
    \item Conclusion: Therefore, there is no fire.
\end{itemize}


$$
p:\text{there is a fire},\qquad q:\text{alarm is ringing}.
$$

$$
\neg q \quad\text{and}\quad \neg p.
$$

Which choice for Premise 1 makes this a correct use of Modus Tollens?

\begin{enumerate}
    \item[(A)] $p\to q$
    \item[(B)] $q\to p$
    \item[(C)] $p$
    \item[(D)] $q$
\end{enumerate}
\end{frame}

\begin{frame}
\begin{block}{Answer: (A)}
Modus Tollens has the form
$$
p\to q,\qquad \neg q
\;\Rightarrow\;
\neg p.
$$

Therefore, Premise 1 must be
$$
p\to q.
$$

Hence,
$$
\boxed{\text{(A)}}
$$
\end{block}
\end{frame}


\begin{frame}{Question 4}
\textbf{Argument:}

\begin{itemize}
    \item Premise 1: \hspace{1cm}
    \item Premise 2: The switch is ON.
    \item Conclusion: Therefore, the light is ON.
\end{itemize}

$$
p:\text{switch is ON},\qquad q:\text{light is ON}.
$$
$$
p\quad\text{and}\quad q.
$$

Which choice for Premise 1 makes this a correct use of Modus Ponens?

\begin{enumerate}
    \item[(A)] $q\to p$
    \item[(B)] $p\to q$
    \item[(C)] $q$
    \item[(D)] $\neg p$
\end{enumerate}
\end{frame}

\begin{frame}
\begin{block}{Answer: (B)}
Modus Ponens has the form
$$
p\to q,\qquad p
\;\Rightarrow\;
q.
$$

Therefore, Premise 1 must be
$$
p\to q.
$$

Hence,
$$
\boxed{\text{(B)}}
$$
\end{block}
\end{frame}


\begin{frame}{Question 5}
\textbf{Which arguments are valid?}
$$
p:\text{battery is dead},\qquad
q:\text{phone will not start}.
$$

\textbf{Argument I:}
$$
p\to q,\qquad q,\qquad\therefore p
$$

\textbf{Argument II:}
$$
p\to q,\qquad p,\qquad\therefore q
$$

\textbf{Argument III:}
$$
p\to q,\qquad \neg q,\qquad\therefore\neg p
$$

\begin{enumerate}
    \item[(A)] I only
    \item[(B)] II only
    \item[(C)] II and III only
    \item[(D)] I, II, and III
\end{enumerate}
\end{frame}

\begin{frame}
\begin{block}{Answer: (C)}
\textbf{I:} 
$$
p\to q,\quad q\;\therefore p
$$
Affirming the consequent $\Rightarrow$ \textbf{invalid}.

\textbf{II:}
$$
p\to q,\quad p\;\therefore q
$$
Modus Ponens $\Rightarrow$ \textbf{valid}.

\textbf{III:}
$$
p\to q,\quad \neg q\;\therefore\neg p
$$
Modus Tollens $\Rightarrow$ \textbf{valid}.

Therefore, the valid arguments are
$$
\boxed{\text{II and III only}}
$$
\end{block}
\end{frame}
%-----------------------------------------------
\section{Tautology and Contradiction}
\begin{frame}{Question 1}

Let $p,q$ be propositional variables.

Consider the formula
$$
F_6:\quad (p\to q)\leftrightarrow(\neg q\to\neg p)
$$

Which of the following is correct?

\begin{enumerate}
    \item[(A)] $F_6$ is a \textbf{tautology}.
    \item[(B)] $F_6$ is a \textbf{contradiction}.
    \item[(C)] $F_6$ is a \textbf{contingency}.
    \item[(D)] $F_6$ is equivalent to $p\land q$.
\end{enumerate}

\end{frame}

\begin{frame}

\begin{block}{Answer: (A)}
$$
\neg q\to\neg p
$$
is the \textbf{contrapositive} of
$$
p\to q.
$$

A statement and its contrapositive always have the same truth value.

Therefore,
$$
(p\to q)\leftrightarrow(\neg q\to\neg p)
$$
is always true.

Hence,
$$
\boxed{\text{(A) }F_6\text{ is a tautology.}}
$$
\end{block}

\end{frame}
\begin{frame}{Question 2}

Let $p,q$ be propositional variables.

Consider the formula
$$
F_7:\quad (p\land q)\land(\neg p\lor\neg q)
$$

Which of the following is correct?

\begin{enumerate}
    \item[(A)] $F_7$ is a \textbf{tautology}.
    \item[(B)] $F_7$ is a \textbf{contradiction}.
    \item[(C)] $F_7$ is a \textbf{contingency}.
    \item[(D)] $F_7$ is equivalent to $p\lor q$.
\end{enumerate}

\end{frame}

\begin{frame}

\begin{block}{Answer: (B)}
$$
p\land q
$$
requires both $p$ and $q$ to be true.

But
$$
\neg p\lor\neg q
$$
requires at least one of $p,q$ to be false.

Therefore, both parts cannot be true at the same time.

Hence,
$$
(p\land q)\land(\neg p\lor\neg q)
$$
is always false.

Thus,
$$
\boxed{\text{(B) }F_7\text{ is a contradiction.}}
$$
\end{block}

\end{frame}
\begin{frame}{Question 3}

\textbf{Domain:} All people.

\textbf{Predicates:}
$$
S(x):\text{``$x$ is a student''},\qquad
H(x):\text{``$x$ is hardworking''}
$$

The statement ``Some students are not hardworking'' is represented as:

\begin{enumerate}
    \item[(A)] $\exists x\,(S(x)\land\neg H(x))$
    \item[(B)] $\exists x\,(S(x)\to\neg H(x))$
    \item[(C)] $\forall x\,(S(x)\to\neg H(x))$
    \item[(D)] $\neg\exists x\,(S(x)\land H(x))$
\end{enumerate}

\end{frame}

\begin{frame}

\begin{block}{Answer: (A)}
``Some students are not hardworking'' means:

There exists a person who is a student and is not hardworking.

$$
\exists x\,(S(x)\land\neg H(x))
$$

Therefore,
$$
\boxed{\text{(A)}}
$$
\end{block}

\end{frame}
\begin{frame}{Question 4}

Let $p$ be a propositional variable.

Consider the formula
$$
F_{10}:\quad (p\to p)\land\neg(p\lor\neg p)
$$

Which of the following is correct?

\begin{enumerate}
    \item[(A)] $F_{10}$ is a \textbf{tautology}.
    \item[(B)] $F_{10}$ is a \textbf{contradiction}.
    \item[(C)] $F_{10}$ is a \textbf{contingency}.
    \item[(D)] $F_{10}$ is equivalent to $p$.
\end{enumerate}

\end{frame}

\begin{frame}

\begin{block}{Answer: (B)}
$$
p\to p
$$
is always true.

Also,
$$
p\lor\neg p
$$
is always true, so
$$
\neg(p\lor\neg p)
$$
is always false.

Therefore,
$$
F_{10}=\text{True}\land\text{False}
=\text{False}.
$$

Hence,
$$
\boxed{\text{(B) }F_{10}\text{ is a contradiction.}}
$$
\end{block}

\end{frame}
\begin{frame}{Question 5}

Let $p,q$ be propositional variables.

The formula
$$
(p\land q)\lor(\neg p\land\neg q)
$$
is:

\begin{enumerate}
    \item[(A)] A tautology
    \item[(B)] A contradiction
    \item[(C)] A contingency
    \item[(D)] Equivalent to $p\lor q$
\end{enumerate}

\end{frame}

\begin{frame}

\begin{block}{Answer: (C)}
The formula is true when $p$ and $q$ have the \textbf{same truth value}.

For $p=T,q=T$:
$$
(T\land T)\lor(F\land F)=T.
$$

For $p=F,q=F$:
$$
(F\land F)\lor(T\land T)=T.
$$

For $p=T,q=F$:
$$
(T\land F)\lor(F\land T)=F.
$$

Thus, the formula is sometimes true and sometimes false.

$$
\boxed{\text{(C) A contingency}}
$$
\end{block}

\end{frame}
\section{Mathematical Induction}
\begin{frame}{Question 1}

Consider the following pseudo-code:

\begin{center}
\texttt{
\begin{tabular}{l}
x := 0\\
y := 0\\
while (n > 0) \{\\
\quad x := x + 1\\
\quad y := y + 2\\
\quad n := n - 1\\
\}
\end{tabular}
}
\end{center}

Initially, $n=N$ for some positive integer.

Which of the following is a correct \textbf{loop invariant}?

\begin{enumerate}
    \item[(A)] $x+y=n$
    \item[(B)] $y=2x$
    \item[(C)] $x=2y$
    \item[(D)] $x=n$
\end{enumerate}

\end{frame}

%------------------------------------------------
\begin{frame}

\begin{block}{Answer: (B)}
Initially,
$$
x=0,\qquad y=0
$$
so
$$
y=2x.
$$

After one iteration,
$$
x'=x+1,\qquad y'=y+2.
$$

If $y=2x$, then
$$
y'=y+2
=2x+2
=2(x+1)
=2x'.
$$

Thus, $y=2x$ remains true after every iteration.

Therefore,
$$
\boxed{\text{(B) }y=2x}
$$
\end{block}

\end{frame}
\begin{frame}{Question 2}

You want to prove a property $P(n)$ for all integers $n\geq2$,
where the proof for $P(k)$ may depend on all previous cases
$P(2),P(3),\ldots,P(k-1)$.

Which proof method is most appropriate?

\begin{enumerate}
    \item[(A)] Simple mathematical induction
    \item[(B)] Strong induction
    \item[(C)] Proof by contradiction only
    \item[(D)] Structural induction
\end{enumerate}

\end{frame}

%------------------------------------------------
\begin{frame}

\begin{block}{Answer: (B)}
In \textbf{strong induction}, we may assume all previous cases
$$
P(2),P(3),\ldots,P(k)
$$
to prove
$$
P(k+1).
$$

Since the proof can depend on all earlier cases, strong induction
is the appropriate method.

$$
\boxed{\text{(B) Strong induction}}
$$
\end{block}

\end{frame}

%------------------------------------------------
\begin{frame}{Question 3}

Arithmetic expressions are built as follows:

\begin{itemize}
    \item Any integer constant $c$ is an expression.
    \item If $E_1,E_2$ are expressions, then $(E_1+E_2)$ is an expression.
    \item If $E_1,E_2$ are expressions, then $(E_1\times E_2)$ is an expression.
\end{itemize}

You want to prove:

\begin{quote}
In every expression, the number of operators is one less than
the number of constants.
\end{quote}

Which correctly describes a structural induction proof?

\begin{enumerate}
    \item[(A)] Base case for a constant; inductive step for
    $(E_1+E_2)$ and $(E_1\times E_2)$.
    \item[(B)] Only prove the property for constants.
    \item[(C)] Assume the property for $(E_1+E_2)$ and prove it for constants.
    \item[(D)] Use simple induction on the numeric value of the constants.
\end{enumerate}

\end{frame}

%------------------------------------------------
\begin{frame}

\begin{block}{Answer: (A)}
\textbf{Base case:} For one constant,
$$
\text{operators}=0,\qquad \text{constants}=1.
$$

Thus,
$$
0=1-1.
$$

\textbf{Inductive step:} Assume the property holds for $E_1$ and
$E_2$. Adding $+$ or $\times$ adds exactly one operator, so the
property remains true.

Hence,
$$
\boxed{\text{(A)}}
$$
\end{block}

\end{frame}

%------------------------------------------------
\begin{frame}{Question 4}

Consider the algorithm:

\begin{center}
\texttt{
\begin{tabular}{l}
result := 1\\
i := 0\\
while (i < N) \{\\
\quad result := result * 2\\
\quad i := i + 1\\
\}
\end{tabular}
}
\end{center}

Which is a suitable invariant to prove that at the end,
$result=2^N$?

\begin{enumerate}
    \item[(A)] At the start of each iteration, $result=2^i$.
    \item[(B)] At the start of each iteration, $result=2^N$.
    \item[(C)] At the start of each iteration, $result=i^2$.
    \item[(D)] At the start of each iteration, $result=2^{N-i}$.
\end{enumerate}

\end{frame}

%------------------------------------------------
\begin{frame}

\begin{block}{Answer: (A)}
Initially,
$$
result=1=2^0=2^i.
$$

After one iteration,
$$
result'=2^i\cdot2=2^{i+1},
\qquad i'=i+1.
$$

Thus,
$$
result'=2^{i'}.
$$

When the loop ends, $i=N$, so
$$
result=2^N.
$$

Therefore,
$$
\boxed{\text{(A) }result=2^i}
$$
is the correct invariant.
\end{block}

\end{frame}
\begin{frame}{Question 5}

Consider the following algorithm for a non-negative integer $N$:

\begin{center}
\texttt{
\begin{tabular}{l}
sum := 0\\
i := 0\\
while (i <= N) \{\\
\quad sum := sum + i\\
\quad i := i + 1\\
\}
\end{tabular}
}
\end{center}

At the end,
$$
sum=0+1+2+\cdots+N=\frac{N(N+1)}{2}.
$$

Which of the following is a suitable \textbf{loop invariant}?

\begin{enumerate}
    \item[(A)] At the start of each iteration, 
    $sum=\frac{i(i+1)}{2}$.
    
    \item[(B)] At the start of each iteration, 
    $sum=\frac{N(N+1)}{2}$.
    
    \item[(C)] At the start of each iteration, 
    $sum=i^2$.
    
    \item[(D)] At the start of each iteration, 
    $sum=\frac{(i-1)i}{2}$.
\end{enumerate}

\end{frame}

%------------------------------------------------
\begin{frame}

\begin{block}{Answer: (D)}
Initially,
$$
i=0,\qquad sum=0.
$$

Thus,
$$
\frac{(i-1)i}{2}
=\frac{(-1)(0)}{2}=0,
$$
so the invariant holds initially.

Assume
$$
sum=\frac{(i-1)i}{2}.
$$

After adding $i$,
$$
sum+i
=
\frac{(i-1)i}{2}+i
=
\frac{i(i+1)}{2}.
$$
\end{block}
\end{frame}

\begin{frame}
    \begin{block}
        
Then $i$ becomes $i+1$. Therefore,
$$
sum=\frac{((i+1)-1)(i+1)}{2}.
$$

Hence the invariant is preserved.

When the loop ends,
$$
i=N+1.
$$

Therefore,
$$
sum=\frac{N(N+1)}{2}.
$$

So the correct answer is
$$
\boxed{\text{(D)}}
$$
\end{block}

\end{frame}
\section{Russell's Paradox}
\begin{frame}{Question 1}

Let $S$ be a set defined by
$$
S=\{\,x\mid x\text{ is a set and }x\notin x\,\}.
$$

Consider the statement: ``$S\in S$''.

Which of the following is correct?

\begin{enumerate}
    \item[(A)] $S\in S$ must be true.
    \item[(B)] $S\in S$ must be false.
    \item[(C)] Assuming $S\in S$ leads to a contradiction, and assuming
    $S\notin S$ also leads to a contradiction.
    \item[(D)] The statement $S\in S$ is always true in standard mathematics.
\end{enumerate}

\end{frame}

%------------------------------------------------

\begin{frame}

\begin{block}{Answer: (C)}
If
$$
S\in S,
$$
then by the definition of $S$,
$$
S\notin S,
$$
which is a contradiction.

On the other hand, if
$$
S\notin S,
$$
then $S$ satisfies the defining condition of $S$, so
$$
S\in S,
$$
which is also a contradiction.

Thus both assumptions lead to contradictions.

$$
\boxed{\text{(C)}}
$$

This is known as \textbf{Russell's Paradox}.
\end{block}

\end{frame}
\begin{frame}{Question 2}

In a certain (naive) set theory, assume:

\begin{quote}
For any property $P(x)$, the collection $\{x\mid P(x)\}$ is a set.
\end{quote}

Let
$$
R=\{x\mid x\notin x\}.
$$

What can be concluded about $R$?

\begin{enumerate}
    \item[(A)] $R$ is a well-defined set with no problems.
    \item[(B)] $R$ is a set, and either $R\in R$ or $R\notin R$, with no contradiction.
    \item[(C)] Treating $R$ as a set leads to a logical contradiction.
    \item[(D)] $R$ must be the empty set.
\end{enumerate}

\end{frame}

\begin{frame}

\begin{block}{Answer: (C)}
By definition,
$$
x\in R\iff x\notin x.
$$

Put $x=R$:
$$
R\in R\iff R\notin R.
$$

Thus, assuming $R$ is a set leads to a logical contradiction.

Therefore,
$$
\boxed{\text{(C)}}
$$
\end{block}

\end{frame}

%------------------------------------------------

\begin{frame}{Question 3}

Suppose we allow the following principle in set theory:

\begin{quote}
For any property $P$, there exists a set of all objects satisfying $P$.
\end{quote}

Let
$$
A=\{x\mid x\text{ is a set and }x\notin x\}.
$$

Which of the following statements are correct?

\begin{enumerate}
    \item[(A)] $A\in A\Rightarrow A\notin A$.
    \item[(B)] $A\notin A\Rightarrow A\in A$.
    \item[(C)] There is no contradiction in assuming $A$ is a set.
    \item[(D)] The existence of $A$ as a set causes a logical inconsistency.
\end{enumerate}

\end{frame}

\begin{frame}

\begin{block}{Answer: (A), (B), (D)}
By definition,
$$
x\in A\iff x\notin x.
$$

Putting $x=A$:

If $A\in A$, then
$$
A\notin A.
$$

If $A\notin A$, then $A$ satisfies the defining property, so
$$
A\in A.
$$

Thus both implications hold, producing a contradiction.

Therefore, (A), (B), and (D) are correct.
\end{block}

\end{frame}

%------------------------------------------------

\begin{frame}{Question 4}

Let
$$
T=\{x\mid x\text{ is a set and }x\notin x\}.
$$

Write the logical equivalence that defines membership in $T$, and then
substitute $x=T$ to show the contradictory situation.

\end{frame}

\begin{frame}

\begin{block}{Answer}
The definition of $T$ gives
$$
\forall x\;(x\in T\iff x\notin x).
$$

Substituting $x=T$:
$$
T\in T\iff T\notin T.
$$

Thus $T$ is an element of itself if and only if it is not an element
of itself, which is a contradiction.
\end{block}

\end{frame}

%------------------------------------------------

\begin{frame}{Question 5}

In a naive set theory where any property defines a set, consider
$$
C=\{x\mid x\text{ is a set and }x\notin x\}.
$$

Which of the following best describes the situation?

\begin{enumerate}
    \item[(A)] $C$ is a perfectly consistent set; either $C\in C$ or
    $C\notin C$, and both are possible in different models.
    
    \item[(B)] $C$ cannot be an element of any set, but can exist as a
    set itself.
    
    \item[(C)] If $C$ is assumed to be a set, then we derive
    $C\in C$ if and only if $C\notin C$.
    
    \item[(D)] $C$ must be equal to the universal set of all sets.
\end{enumerate}

\end{frame}

\begin{frame}

\begin{block}{Answer: (C)}
By definition,
$$
x\in C\iff x\notin x.
$$

Putting $x=C$ gives
$$
C\in C\iff C\notin C.
$$

Therefore, assuming $C$ is a set leads directly to a contradictory
equivalence.

Hence,
$$
\boxed{\text{(C)}}
$$
\end{block}

\end{frame}
\section{Proofs}
%------------------------------------------------
\begin{frame}{Question 1}

Let $a,b,c$ be integers. Consider:

\begin{quote}
If $a\mid b$ and $a\mid c$, then $a\mid(b+c)$.
\end{quote}

Which proof pattern is used?

\begin{enumerate}
    \item[(A)] Assume $a\mid(b+c)$ and prove $a\mid b$ and $a\mid c$.
    \item[(B)] Assume $a\mid b$ and $a\mid c$, then prove $a\mid(b+c)$.
    \item[(C)] Assume $a\nmid(b+c)$ and prove $a\nmid b$ or $a\nmid c$.
    \item[(D)] Assume $a\mid b$ or $a\mid c$, then prove $a\mid(b+c)$.
\end{enumerate}

\end{frame}

%------------------------------------------------
\begin{frame}

\begin{block}{Answer: (B)}
To prove
$$
P\land Q\Rightarrow R,
$$
we assume $P$ and $Q$ and then prove $R$.

Here,
$$
a\mid b,\qquad a\mid c
$$
are assumed, and we prove
$$
a\mid(b+c).
$$

Therefore, the correct proof pattern is

$$
\boxed{\text{(B)}}
$$
\end{block}

\end{frame}

%------------------------------------------------
\begin{frame}{Question 2}

To prove: For any real $x$,
$$
|x|^2=x^2.
$$

Which proof method is most natural?

\begin{enumerate}
    \item[(A)] Proof by contradiction
    \item[(B)] Proof by cases
    \item[(C)] Strong induction
    \item[(D)] Structural induction
\end{enumerate}

\end{frame}

%------------------------------------------------
\begin{frame}

\begin{block}{Answer: (B)}
Consider two cases.

\textbf{Case 1:} $x\geq0$.
Then
$$
|x|=x\quad\Rightarrow\quad |x|^2=x^2.
$$

\textbf{Case 2:} $x<0$.
Then
$$
|x|=-x\quad\Rightarrow\quad |x|^2=(-x)^2=x^2.
$$

Thus the statement is proved by considering the two possible cases.

$$
\boxed{\text{(B) Proof by cases}}
$$
\end{block}

\end{frame}

%------------------------------------------------
\begin{frame}{Question 3}

Which of the following is best proved by \textbf{contradiction}?

\begin{enumerate}
    \item[(A)] If $n$ is even, then $n+2$ is even.
    \item[(B)] There is no largest integer.
    \item[(C)] $A\cup B=B\cup A$ for all sets $A,B$.
    \item[(D)] If $a=b$, then $a-b=0$.
\end{enumerate}

\end{frame}

%------------------------------------------------
\begin{frame}

\begin{block}{Answer: (B)}
Assume, for contradiction, that there is a largest integer $M$.

But
$$
M+1>M,
$$
so $M+1$ is an integer larger than $M$.

This contradicts the assumption that $M$ was the largest integer.

Therefore, there is no largest integer.

$$
\boxed{\text{(B)}}
$$
\end{block}

\end{frame}

%------------------------------------------------
\begin{frame}{Question 4}

For any sets $A,B,C$, which identities are always true?

\begin{enumerate}
    \item[(A)] $A\cap(B\cup C)=(A\cap B)\cup(A\cap C)$
    \item[(B)] $A\cup(B\cap C)=(A\cup B)\cap(A\cup C)$
    \item[(C)] $A\setminus(B\cup C)=(A\setminus B)\cup(A\setminus C)$
    \item[(D)] $A\setminus(B\cap C)=(A\setminus B)\cap(A\setminus C)$
\end{enumerate}

\end{frame}

%------------------------------------------------
\begin{frame}

\begin{block}{Answer: (A), (B)}
(A) and (B) are the \textbf{distributive laws} for sets:
$$
A\cap(B\cup C)=(A\cap B)\cup(A\cap C),
$$
$$
A\cup(B\cap C)=(A\cup B)\cap(A\cup C).
$$

(C) is false. The correct identity is
$$
A\setminus(B\cup C)
=(A\setminus B)\cap(A\setminus C).
$$

(D) is false. The correct identity is
$$
A\setminus(B\cap C)
=(A\setminus B)\cup(A\setminus C).
$$

Therefore,
$$
\boxed{\text{(A), (B)}}
$$
\end{block}

\end{frame}
%------------------------------------------------
\begin{frame}{Question 5}

To prove:

\begin{quote}
For integers $m,n$, $m+n=n+m$ if and only if addition is
commutative for integers.
\end{quote}

Which of the following must be shown?

\begin{enumerate}
    \item[(A)] If $m+n=n+m$, then addition is commutative for integers.
    \item[(B)] If addition is commutative for integers, then $m+n=n+m$.
    \item[(C)] If $m+n\neq n+m$, then addition is not commutative.
    \item[(D)] Both directions:
    \begin{itemize}
        \item[(i)] $m+n=n+m\Rightarrow$ commutativity,
        \item[(ii)] commutativity $\Rightarrow m+n=n+m$.
    \end{itemize}
\end{enumerate}

\end{frame}

%------------------------------------------------
\begin{frame}

\begin{block}{Answer: (D)}
An ``if and only if'' statement
$$
P\iff Q
$$
requires proving both directions:
$$
P\Rightarrow Q
\qquad\text{and}\qquad
Q\Rightarrow P.
$$

Here,
$$
P:m+n=n+m
$$
and
$$
Q:\text{addition is commutative}.
$$

Therefore, both directions must be shown.

Option (B) gives only one direction, while (D) gives both.

$$
\boxed{\text{(D)}}
$$
\end{block}

\end{frame}
\end{document}